\chapter{Hito 7: Construcción de la Capa Core y Dimensiones en DuckDB}

\section{Objetivo}

El objetivo de este hito es construir la capa \texttt{core} del Data Warehouse.

En esta capa comenzaremos a transformar los datos crudos almacenados en \texttt{staging} en estructuras analíticas organizadas siguiendo principios de modelado dimensional.

Al finalizar este hito el estudiante será capaz de:

\begin{itemize}
\item Crear dimensiones en DuckDB.
\item Comprender el uso de claves subrogadas.
\item Construir una dimensión tiempo.
\item Construir una dimensión moneda.
\item Construir una dimensión entidad financiera.
\item Poblar dimensiones a partir de los datos de staging.
\item Preparar la base para construir las tablas de hechos.
\end{itemize}

\section{Resultado esperado}

Al finalizar este hito deberán existir:

\begin{verbatim}
core.dim_tiempo

core.dim_moneda

core.dim_entidad_financiera
\end{verbatim}

Estas dimensiones serán utilizadas posteriormente por las tablas de hechos.

\section{Arquitectura del hito}

El flujo de transformación será:

\begin{verbatim}
staging
|
V
core.dim_tiempo

core.dim_moneda

core.dim_entidad_financiera
\end{verbatim}

\section{Paso 1: Crear script de dimensiones}

Crear:

\begin{verbatim}
sql/03_dimensions_ddl.sql
\end{verbatim}

\section{Paso 2: Crear dimensión tiempo}

Agregar:

\begin{lstlisting}[language=SQL]
CREATE TABLE IF NOT EXISTS core.dim_tiempo
(
sk_tiempo INTEGER PRIMARY KEY,

```
fecha DATE NOT NULL,

anio INTEGER,

mes INTEGER,

trimestre INTEGER,

anio_mes INTEGER,

nombre_mes VARCHAR,

es_fin_de_mes BOOLEAN
```

);
\end{lstlisting}

\section{Paso 3: Crear dimensión moneda}

Agregar:

\begin{lstlisting}[language=SQL]
CREATE TABLE IF NOT EXISTS core.dim_moneda
(
sk_moneda INTEGER PRIMARY KEY,

```
moneda_cod VARCHAR,

moneda_nombre VARCHAR
```

);
\end{lstlisting}

\section{Paso 4: Crear dimensión entidad financiera}

Agregar:

\begin{lstlisting}[language=SQL]
CREATE TABLE IF NOT EXISTS
core.dim_entidad_financiera
(
sk_entidad INTEGER PRIMARY KEY,

```
entidad_cod VARCHAR,

entidad_nombre VARCHAR,

tipo_entidad VARCHAR,

vigente_desde DATE,

vigente_hasta DATE,

es_actual BOOLEAN
```

);
\end{lstlisting}

\section{Paso 5: Ejecutar DDL}

Desde Python:

\begin{lstlisting}[language=Python]
from db import get_connection

conn = get_connection()

with open(
"sql/03_dimensions_ddl.sql",
"r",
encoding="utf-8"
) as file:

```
conn.execute(
    file.read()
)
```

conn.close()
\end{lstlisting}

\section{Paso 6: Crear script de carga}

Crear:

\begin{verbatim}
sql/04_load_dimensions.sql
\end{verbatim}

\section{Paso 7: Limpiar dimensiones}

Agregar:

\begin{lstlisting}[language=SQL]
DELETE FROM
core.dim_entidad_financiera;

DELETE FROM
core.dim_moneda;

DELETE FROM
core.dim_tiempo;
\end{lstlisting}

\section{Paso 8: Poblar dimensión tiempo}

Agregar:

\begin{lstlisting}[language=SQL]
INSERT INTO core.dim_tiempo

SELECT

```
CAST(
    strftime(
        fecha,
        '%Y%m%d'
    ) AS INTEGER
) AS sk_tiempo,

fecha,

year(fecha) AS anio,

month(fecha) AS mes,

quarter(fecha) AS trimestre,

CAST(
    strftime(
        fecha,
        '%Y%m'
    ) AS INTEGER
) AS anio_mes,

monthname(fecha)
    AS nombre_mes,

fecha =
last_day(fecha)
    AS es_fin_de_mes
```

FROM generate_series(
DATE '2022-01-01',
DATE '2024-12-31',
INTERVAL 1 DAY
) t(fecha);
\end{lstlisting}

\section{Explicación}

DuckDB posee una función muy útil:

\begin{lstlisting}[language=SQL]
generate_series()
\end{lstlisting}

que permite generar calendarios completos sin necesidad de archivos externos.

\section{Paso 9: Poblar dimensión moneda}

Agregar:

\begin{lstlisting}[language=SQL]
INSERT INTO core.dim_moneda

SELECT

```
ROW_NUMBER()
OVER(
    ORDER BY moneda
) AS sk_moneda,

moneda AS moneda_cod,

CASE

    WHEN moneda='USD'
    THEN 'Dolar estadounidense'

    WHEN moneda='EUR'
    THEN 'Euro'

    ELSE moneda

END AS moneda_nombre
```

FROM
(
SELECT DISTINCT
moneda
FROM staging.tipo_cambio_raw
) x;
\end{lstlisting}

\section{Paso 10: Poblar dimensión entidad financiera}

Agregar:

\begin{lstlisting}[language=SQL]
INSERT INTO
core.dim_entidad_financiera

SELECT

```
ROW_NUMBER()
OVER(
    ORDER BY entidad_cod
) AS sk_entidad,

entidad_cod,

entidad_nombre,

tipo_entidad,

DATE '2022-01-01'
    AS vigente_desde,

DATE '9999-12-31'
    AS vigente_hasta,

TRUE
    AS es_actual
```

FROM
(
SELECT DISTINCT

```
    entidad_cod,

    entidad_nombre,

    tipo_entidad

FROM
    staging.balance_bancario_raw
```

) x;
\end{lstlisting}

\section{Nota sobre SCD Tipo 2}

En esta primera versión no implementaremos todavía una dimensión histórica completa.

La lógica será:

\begin{itemize}
\item Una fila por entidad.
\item Una sola versión vigente.
\item Sin cambios históricos.
\end{itemize}

En una versión posterior agregaremos:

\begin{itemize}
\item Versionamiento.
\item Fechas de vigencia.
\item SCD Tipo 2 completo.
\end{itemize}

\section{Paso 11: Ejecutar carga}

Crear:

\begin{verbatim}
etl/load_dimensions.py
\end{verbatim}

Contenido:

\begin{lstlisting}[language=Python]
from db import get_connection

conn = get_connection()

with open(
"sql/04_load_dimensions.sql",
"r",
encoding="utf-8"
) as file:

```
sql = file.read()
```

conn.execute(sql)

conn.close()

print(
"Dimensiones cargadas."
)
\end{lstlisting}

\section{Paso 12: Verificar dimensión tiempo}

Crear:

\begin{verbatim}
sql/check_dim_tiempo.sql
\end{verbatim}

\begin{lstlisting}[language=SQL]
SELECT *
FROM core.dim_tiempo
LIMIT 10;
\end{lstlisting}

\section{Paso 13: Verificar dimensión moneda}

\begin{lstlisting}[language=SQL]
SELECT *
FROM core.dim_moneda;
\end{lstlisting}

Resultado esperado:

\begin{verbatim}
USD
EUR
\end{verbatim}

\section{Paso 14: Verificar dimensión entidad financiera}

\begin{lstlisting}[language=SQL]
SELECT *
FROM core.dim_entidad_financiera;
\end{lstlisting}

\section{Paso 15: Crear consulta de monitoreo}

Crear:

\begin{verbatim}
sql/check_dimensions.sql
\end{verbatim}

Contenido:

\begin{lstlisting}[language=SQL]
SELECT

```
'dim_tiempo' AS tabla,

COUNT(*) AS total_filas
```

FROM core.dim_tiempo

UNION ALL

SELECT

```
'dim_moneda',

COUNT(*)
```

FROM core.dim_moneda

UNION ALL

SELECT

```
'dim_entidad_financiera',

COUNT(*)
```

FROM core.dim_entidad_financiera;
\end{lstlisting}

\section{Paso 16: Ejecutar monitoreo}

Resultado esperado:

\begin{verbatim}
dim_tiempo

dim_moneda

dim_entidad_financiera
\end{verbatim}

con conteos positivos.

\section{Paso 17: Integrar al pipeline}

Actualizar:

\begin{verbatim}
etl/run_pipeline.py
\end{verbatim}

\begin{lstlisting}[language=Python]
import subprocess

subprocess.run(
[
"python",
"etl/load_staging.py"
],
check=True
)

subprocess.run(
[
"python",
"etl/load_dimensions.py"
],
check=True
)

print(
"Pipeline ejecutado correctamente."
)
\end{lstlisting}

\section{Buenas prácticas}

Durante este hito seguiremos las siguientes reglas:

\begin{enumerate}
\item Las dimensiones siempre se construyen desde staging.
\item Las claves subrogadas pertenecen al Data Warehouse.
\item La dimensión tiempo debe ser reutilizable.
\item La dimensión entidad debe ser preparada para soportar SCD.
\item Las cargas deben ser reproducibles.
\end{enumerate}

\section{Checklist de validación}

\begin{itemize}
\item[$\Box$] Dimensión tiempo creada.
\item[$\Box$] Dimensión moneda creada.
\item[$\Box$] Dimensión entidad creada.
\item[$\Box$] Dimensión tiempo poblada.
\item[$\Box$] Dimensión moneda poblada.
\item[$\Box$] Dimensión entidad poblada.
\item[$\Box$] Script de monitoreo creado.
\item[$\Box$] Pipeline actualizado.
\end{itemize}

\section{Entregable}

Al finalizar este hito el estudiante deberá disponer de:

\begin{itemize}
\item Tres dimensiones pobladas.
\item Scripts SQL documentados.
\item Script de carga automatizado.
\item Pipeline funcional.
\item Capa Core parcialmente construida.
\end{itemize}

\section{Próximo hito}

En el Hito 8 construiremos las tablas de hechos.

Aprenderemos a:

\begin{itemize}
\item Resolver claves subrogadas.
\item Crear fact_tipo_cambio.
\item Crear fact_balance_mensual.
\item Calcular medidas derivadas.
\item Validar el grano de cada hecho.
\end{itemize}
